\section{Introduction}
\label{sec:introduction}
Investment has been seen as an important part of explaining the business cycle for a long time. It remains the most volatile variable in the macroeconomy. Over the past few decades, the role of investment in the business cycle has been continuously discussed, with new model environments being proposed \citep{greenwoodInvestmentCapacityUtilization1988,greenwoodLongRunImplicationsInvestmentSpecific1997a,fisherDynamicEffectsNeutral2006,justinianoInvestmentShocksBusiness2010,winberryLumpyInvestmentBusiness2021}.

Early research has drawn attention to the investment shock because of the obvious trend in the data on the relative price between investment goods and consumption goods, which implies another trend term in the investment specific technology besides the neutral technology \citep{greenwoodInvestmentCapacityUtilization1988}. After that, quantitative results have been established by \cite{greenwoodLongRunImplicationsInvestmentSpecific1997a} using a calibrated DSGE model, and by \cite{fisherDynamicEffectsNeutral2006} using a structural VAR model. Both of them claimed that the investment shock accounts for about half of the forecast error of the output and the working hours.

A new wave of research about the investment shock is stimulated by \cite{smetsShocksFrictionsUS2007}, who estimated a more frictional model based on \cite{christianoNominalRigiditiesDynamic2005} using Bayesian estimation. Using the same approach, \cite{justinianoInvestmentShocksBusiness2010, justinianoInvestmentShocksRelative2011} argued again that the investment shocks are the main source of the business cycle. While the investment shock shows strong effects on the macroeconomic variables, there are other candidates being proposed to account for the business cycle. \cite{christianoRiskShocks2014} and \cite{kamberFinancialFrictionsRole2015} argued that investment shock is less important once the model includes financial friction and shocks in the financial market.

Though I use the name of investment shock in the previous paragraphs, there are various kinds of specification being discussed. Among them, the most classic one is the investment specific technology (IST) shock, which increases the technology level to produce the investment goods by the consumption goods \citep{greenwoodLongRunImplicationsInvestmentSpecific1997a, fisherDynamicEffectsNeutral2006, justinianoInvestmentShocksBusiness2010}. Besides, \cite{justinianoInvestmentShocksRelative2011, kamberFinancialFrictionsRole2015} distinguished the IST shock from the marginal efficiency of investment (MEI) shock, which represents the technology level to produce the capital goods by the investment goods. However, the difference between these two shocks rely on the detailed modeling of the capital production process. If the capital production is modelled by one capital law of motion equation, then those two shocks cannot be distinguished in their effect, and thus I make no clarification between them except for Section~\ref{sec:robustness:mei}. In addition, there are many works discussed the IST news shock, in the sense that the shock is known by agents now but will affect the economy several periods after \citep{schmitt-groheWhatsNewsBusiness2012, benzeevInvestmentSpecificNewsShocks2015, liaoNewsShocksInvestmentspecific2023}. However, this kind of shock is beyond the scope of this paper.

Despite the IST shock has been considered as an important driving force of the business cycle, it fails to generate the comovement between the consumption, investment and inflation. To solve the comovement problem, \cite{nireiOriginsPropagationAnimal} proposed an investment demand shock based on a model featuring the firms' heterogeneity and the lumpy investment assumption. By the calibrated model, the investment demand shock generates the comovement between consumption and investment. While the shock is traced back to the idiosyncratic productivity shock, on the aggregate level, in macro level it presents a deviation of investment from the original investment decision made in the shock-free environment. Because of this, the investment demand shock can also be represented in a DSGE model without heterogeneity, and be estimated using the same approach with \cite{smetsShocksFrictionsUS2007}.

However, since this shock requires to fix the endogenous investment decision, the solution method of the rational expectation equilibrium needs to be separated into two steps. In the first step, the model includes all shocks and frictions except for the investment demand shock, from where the endogenous investment decision is solved. In the second step, the investment is decided by both the endogenous investment decision and the exogenous investment demand shock. It is worth pointing out that the endogenous investment decision cannot be changed in the second step, and thus investment demand is exogenously increased by the shock. Although the investment decision is fixed, households can still adjust their decisions on other terms, such as consumption and labor input, and thus the second step solves another equilibrium. This equilibrium, finally, will be used to compute the likelihood.

The identification strategy is not direct---it is hard to directly identify invesetment demand shock by data, as the shock is originated from idiosyncratic productivity shock. Moreover, though we can use relative price of invesetment goods as the proxy for IST shock in some sense, there is still a concern that we might attribute all residuals to invesetment demand shock. Therefore, I postpone this exercise to a later section and rely on the theoretical difference---whether the comovement between consumption, investment and inflation---to identify two shocks in the baseline result. However, due to the high-dimensional parameter space, the comovement difference is hard to test among all combinations of parameter values. Therefore, I claim the shortage of such identification strategy.

Section~\ref{sec:model} introduces the baseline model. Section~\ref{sec:solution} presents the solution method and incorporates the investment demand shock. Section~\ref{sec:estimation} describes the Bayesian estimation, and Section~\ref{sec:results} reports the benchmark results. Section~\ref{sec:robustness} provides robustness checks, and Section~\ref{sec:conclusion} concludes.